\chapter{Parathyroidectomy}
\section*{Overview}
Parathyroidectomy removes one or more overactive parathyroid glands to correct excessive parathyroid hormone production.
\section*{Common Instruments and Equipment}
The exact setup varies by technique and setting. Common items may include:
\begin{itemize}[nosep,leftmargin=*]
  \item neck surgery instrument set and magnification
  \item intraoperative parathyroid hormone testing supplies when used
  \item nerve-monitoring equipment when indicated
\end{itemize}

\section*{Why It Is Done}
It is commonly performed for primary hyperparathyroidism with significant hypercalcemia, kidney stones, bone disease, symptoms, or selected high-risk features.
\section*{How It Is Performed}
Through a small neck incision, the surgeon identifies and removes the abnormal gland or glands. Intraoperative hormone testing and focused or bilateral exploration may be used.
\section*{Contraindications}
No single contraindication applies to every patient; the treating team assesses whether the procedure is suitable for the individual and the reason for treatment.

\section*{Precautions}
The care team reviews relevant health conditions, medicines, allergies, and procedure-specific risks, and explains any changes needed for the individual.

\section*{Risks and Recovery}
Low calcium, bleeding, recurrent laryngeal-nerve injury, persistent disease, and recurrent hyperparathyroidism are possible. Calcium monitoring and supplementation may be needed after surgery.
\section*{References}
\begin{enumerate}\item MedlinePlus. Parathyroidectomy. U.S. National Library of Medicine; 2025.\item Current Medical Diagnosis and Treatment. McGraw-Hill; 2025.\end{enumerate}
\clearpage
